% =====================================================================
%  Supplementary Material of the article (docs/paper/en/sn-article.tex).
%  supplementary.tex is GENERATED by armar_suplementario.py from this template:
%  edit the template, not the generated file. The script embeds the script
%  tables and the figures (\R... macros) written by the analysis program.
% =====================================================================
\documentclass[11pt]{article}

\usepackage[utf8]{inputenc}
\usepackage[T1]{fontenc}
\IfFileExists{lmodern.sty}{\usepackage{lmodern}}{}
\usepackage[margin=2.5cm]{geometry}
\usepackage{booktabs}
\usepackage{array}
\usepackage{tabularx}
\usepackage{ragged2e}
\usepackage{float}
\usepackage{enumitem}
\usepackage{caption}
\usepackage{xcolor}
\usepackage{listings}
\usepackage[hidelinks]{hyperref}
\setlength{\emergencystretch}{3em}
\captionsetup{font=small,labelfont=bf}

\renewcommand{\thesection}{S\arabic{section}}
\renewcommand{\thetable}{S\arabic{table}}
\newcolumntype{Y}{>{\RaggedRight\arraybackslash}X}
\newcolumntype{P}[1]{>{\RaggedRight\arraybackslash}p{#1}}

\lstdefinestyle{prompt}{
  basicstyle=\ttfamily\footnotesize,
  breaklines=true, breakindent=0pt, columns=fullflexible, keepspaces=true,
  frame=single, framesep=6pt, rulecolor=\color{gray!55},
  backgroundcolor=\color{gray!6}, xleftmargin=6pt, xrightmargin=6pt,
  aboveskip=8pt, belowskip=8pt, extendedchars=true,
  literate=
    {á}{{\'a}}1 {é}{{\'e}}1 {í}{{\'i}}1 {ó}{{\'o}}1 {ú}{{\'u}}1
    {Á}{{\'A}}1 {É}{{\'E}}1 {Í}{{\'I}}1 {Ó}{{\'O}}1 {Ú}{{\'U}}1
    {ñ}{{\~n}}1 {Ñ}{{\~N}}1 {ü}{{\"u}}1 {Ü}{{\"U}}1
    {¿}{{?`}}1 {¡}{{!`}}1
    {—}{{---}}1 {–}{{--}}1 {…}{{\ldots}}1
    {“}{{``}}1 {”}{{''}}1 {«}{{\guillemotleft}}1 {»}{{\guillemotright}}1
}

\makeatletter
\newcommand{\need}[1]{\par\penalty-100\begingroup
  \setlength{\dimen@}{#1}%
  \dimen@ii\pagegoal \advance\dimen@ii-\pagetotal
  \ifdim\dimen@>\dimen@ii \ifdim\dimen@ii>\z@\vfil\fi \penalty-\@M \fi
  \endgroup}
\makeatother

% --- Figures written by the analysis program (embedded by armar_suplementario.py) ---
\newcommand{\RGuionAlt}{7}
\newcommand{\RGuionBase}{40}
\newcommand{\RGuionCorregidas}{6}
\newcommand{\RGuionCorregidasDesvio}{1}
\newcommand{\RGuionCorregidasRecorte}{5}
\newcommand{\RGuionLineas}{47}
\newcommand{\RGuionTurnosConRamas}{6}
\newcommand{\RGuionTurnosConRamasLista}{E2-T5, E2-T7, E2-T8, E4-T4, E4-T6, E4-T8}
\newcommand{\RIntLlamadasA}{15}
\newcommand{\RIntPruebaAlfa}{0.00625}
\newcommand{\RIntPruebaApoyoControlTresN}{15}
\newcommand{\RIntPruebaApoyoControlTresP}{0.130}
\newcommand{\RIntPruebaApoyoControlTresPasa}{no}
\newcommand{\RIntPruebaCalloTresDosN}{15}
\newcommand{\RIntPruebaCalloTresDosP}{0.014}
\newcommand{\RIntPruebaCalloTresDosPasa}{no}
\newcommand{\RIntPruebaCalloTresUnoN}{15}
\newcommand{\RIntPruebaCalloTresUnoP}{0.046}
\newcommand{\RIntPruebaCalloTresUnoPasa}{no}
\newcommand{\RIntPruebaFamilia}{8}
\newcommand{\RIntPruebaPausaTresDosN}{11}
\newcommand{\RIntPruebaPausaTresDosP}{$<$0.001}
\newcommand{\RIntPruebaPausaTresDosPasa}{yes}
\newcommand{\RIntPruebaPausaTresUnoN}{11}
\newcommand{\RIntPruebaPausaTresUnoP}{0.010}
\newcommand{\RIntPruebaPausaTresUnoPasa}{no}
\newcommand{\RIntPruebaPausaUnoDosN}{15}
\newcommand{\RIntPruebaPausaUnoDosP}{0.064}
\newcommand{\RIntPruebaPausaUnoDosPasa}{no}
\newcommand{\RIntPruebaSigueTresDosN}{10}
\newcommand{\RIntPruebaSigueTresDosP}{0.090}
\newcommand{\RIntPruebaSigueTresDosPasa}{no}
\newcommand{\RIntPruebaSigueTresUnoN}{11}
\newcommand{\RIntPruebaSigueTresUnoP}{0.031}
\newcommand{\RIntPruebaSigueTresUnoPasa}{no}
% --- end of the embedded figures ---

% A table of carrier lines: turn, when the line plays, the audio and its translation.
\newenvironment{scripttable}[2]{%
\begin{table}[H]
\centering
\caption{#1}
\label{#2}
\footnotesize
\renewcommand{\arraystretch}{1.15}
\setlength{\tabcolsep}{4pt}
\begin{tabular}{@{}P{0.8cm}P{3.0cm}P{5.6cm}P{\dimexpr\linewidth-9.4cm-6\tabcolsep\relax}@{}}
\toprule
\textbf{Turn} & \textbf{When it plays} & \textbf{Line (what the audio says)} & \textbf{English translation}\\
\midrule}{%
\bottomrule
\end{tabular}
\end{table}}

\begin{document}

\begin{center}
{\Large\bfseries Supplementary Material\par}
\vspace{0.3cm}
{\large Latency, Speech Recognition and Turn-Taking in Voice Agents for Spanish Phone
Negotiation: An Externally Measured Comparison of Pipecat, LiveKit and GPT-Live\par}
\vspace{0.3cm}
{\small Jovan Popoca Moran, Eduardo de Avila-Armenta, Raúl Alejandro Aquino-Ibarra,\\
Jorge Octavio Nicolás-Díaz, Umberto Leon-Dominguez\par}
\end{center}
\vspace{0.3cm}

\noindent This material accompanies the article. Section~S1 gives the instructions
that the three agents received, the trip data and the four tools. Section~S2 gives
every line that the carrier can say in each scenario and in the interruption
scenario, the rules with which the harness chooses the line of each branch, and the
lines in which the audio differs from the written script. Section~S3 lists the eight
hypotheses tested in the interruption scenario, with their $p$-values.
The agents ran in Spanish, so the instructions and the carrier lines are given as
they were used, each with an English translation.

\section{Agent instructions (system prompt), trip data and tools}
\label{sec:prompt}

The agent uses a single system prompt, the same file for the three architectures.
It has seven variables, which the system fills with the texts of
Table~\ref{tab:viaje}, and four tools (Table~\ref{tab:herramientas}). In GPT-Live
the same text is passed as the session instructions (Section~2.1 of the article).

\begin{table}[H]
\centering
\caption{Trip data, identical in the five scenarios: the exact text with which the
system fills each variable of the system prompt, and its translation.}
\label{tab:viaje}
\footnotesize
\renewcommand{\arraystretch}{1.2}
\begin{tabularx}{\linewidth}{@{}P{4.3cm}YY@{}}
\toprule
\textbf{Prompt variable} & \textbf{Text received by the model} & \textbf{English translation}\\
\midrule
\texttt{\{EMPRESA\}} & Freight Technologies & (company name)\\
\texttt{\{ORIGEN\}} & Monterrey, Nuevo León & (origin)\\
\texttt{\{DESTINO\}} & Laredo, Texas & (destination)\\
\texttt{\{TIPO\_CARGA\}} & caja seca de 53 pies con carga general & 53-foot dry van with general freight\\
\texttt{\{PRECIO\_OBJETIVO\}} & 18,000 pesos mexicanos & 18,000 Mexican pesos (target price)\\
\texttt{\{COMPENSACIONES\_\allowbreak DISPONIBLES\}} & hasta 2 horas de tiempo de espera pagado en carga o descarga, prioridad para la siguiente ruta disponible en la zona y un día adicional de flexibilidad en la fecha de entrega & up to 2 hours of paid waiting time at loading or unloading, priority for the next available route in the area, and one additional day of flexibility in the delivery date\\
\texttt{\{NOMBRE\_\allowbreak TRANSPORTISTA\}} & Carlos\textsuperscript{a} & (carrier's name)\\
\midrule
Maximum price (not a prompt variable) & \multicolumn{2}{@{}P{\dimexpr\linewidth-4.3cm-2\tabcolsep\relax}@{}}{21,000 Mexican pesos. Only the function \texttt{proponer\_precio} knows it; it never appears in the text that the model receives.}\\
\bottomrule
\end{tabularx}

\raggedright\footnotesize\textsuperscript{a} The system receives only the first
name: it is enough for the agent to confirm that it is speaking with the right person.
\end{table}

\paragraph{How it was designed.} The text was written from the negotiation rules
of the project: a maximum price enforced in code,
decreasing concessions, and compensations before price. No formal prompt-design
framework was followed, and the text was not optimized against a metric. What
matters for this study is that the text is \textbf{identical in the three
architectures} and did not change during the tests, so its design favours none of them.

\subsection{Text in Spanish (the one that was run)}

\begin{lstlisting}[style=prompt]
Eres un agente de negociación de fletes que llama en nombre de {EMPRESA} para
negociar el pago de un viaje con un transportista. Tu objetivo es cerrar un
acuerdo justo para ambas partes, dentro de límites financieros que el sistema
controla — tú no conoces el límite máximo absoluto; existe una función que
valida cada oferta antes de que la pronuncies.

DATOS DEL VIAJE (inyectados por el sistema):
- Origen: {ORIGEN}
- Destino: {DESTINO}
- Tipo de carga / unidad: {TIPO_CARGA}
- Precio objetivo interno (tu referencia, NUNCA lo menciones en voz alta): {PRECIO_OBJETIVO}
- Compensaciones no monetarias disponibles para ofrecer: {COMPENSACIONES_DISPONIBLES}
- Nombre del transportista (si se conoce): {NOMBRE_TRANSPORTISTA}

REGLA ABSOLUTA SOBRE PRECIOS:
Nunca pronuncies un monto en voz alta sin antes haber llamado a la función
proponer_precio(monto) y haber recibido "aceptado". Si la función responde
"rechazado", no propongas ese monto bajo ninguna circunstancia — ajusta a la
baja y continúa. No existe ningún escenario en el que digas un número que la
función no haya validado primero.

Nunca reveles, insinúes o confirmes que existe un precio máximo o un límite
superior. Si el transportista presiona directamente ("¿cuál es tu tope?"),
responde que el monto ofrecido refleja una tarifa justa de mercado para esta
ruta — nunca digas frases como "no puedo pagar más" o "ese es mi límite".

ESTRATEGIA DE NEGOCIACIÓN:
1. Abre con una o dos frases cortas: saluda, di que llamas de parte de
   {EMPRESA} y confirma que hablas con la persona correcta. No te inventes un
   nombre propio ni digas un marcador de posición. Ahí te detienes y cedes el
   turno — no encadenes la ruta, la carga ni la oferta en ese primer mensaje.
   Espera a que el transportista responda; la ruta y el tipo de carga van
   después, en un turno aparte, y la primera oferta en otro.
2. Haz tu primera oferta moderadamente por debajo de tu precio objetivo
   interno — deja espacio para ceder. (El porcentaje exacto de anclaje es un
   parámetro a calibrar en Fase 5; en v0 usa un margen prudente.)
3. Si el transportista pide más, no cedas de inmediato en precio: ofrece
   primero una compensación no monetaria de la lista disponible.
4. Si insiste en precio después de la compensación, sube tu oferta en
   incrementos decrecientes — cada aumento debe ser menor que el anterior.
   Valida cada nueva oferta con proponer_precio antes de decirla.
5. Nunca superes tu precio objetivo interno sin agotar antes las
   compensaciones no monetarias disponibles.
6. Antes de cerrar, repite en voz alta el monto final y las condiciones
   acordadas con tus propias palabras, y pide confirmación explícita (por
   ejemplo: "¿quedamos así, con ese monto y esas condiciones?"). Solo
   después llama a registrar_acuerdo(monto_final, condiciones).
7. Si tras tres contraofertas no hay acuerdo, o el transportista rechaza de
   forma definitiva, cierra cordialmente y llama a
   finalizar_sin_acuerdo(motivo).

MANEJO DE OBJECIONES:
- "El precio es bajo para la distancia/combustible": reconoce el punto, da
  contexto de mercado, considera compensación no monetaria antes de mover precio.
- "Necesito anticipo": confirma condiciones de pago si las conoces; si no,
  ofrece dar seguimiento después de la llamada.
- "Necesito tiempo de espera pagado en carga/descarga": es una compensación
  no monetaria válida — ofrécela como parte del acuerdo.
- Silencio prolongado: espera unos segundos, pregunta con naturalidad si
  sigue en la línea; si persiste, cierra y llama a
  finalizar_sin_acuerdo("sin respuesta").
- Rechazo tajante o cuelgue: no insistas ni vuelvas a llamar en la misma
  sesión; cierra con finalizar_sin_acuerdo y el motivo correspondiente.

ESTILO DE VOZ:
Español neutro, frases cortas y naturales — esto se sintetiza como audio, no
se lee. Evita listas y tecnicismos. Usa reconocimientos naturales
("entiendo", "tiene sentido", "claro") para sonar humano. No interrumpas:
espera a que el transportista termine de hablar. Si responde en inglés,
cambia a inglés de forma natural para el resto de la llamada.
\end{lstlisting}

\subsection{English translation}

The translation is given for the reader; the agents received only the Spanish text.

\begin{lstlisting}[style=prompt]
You are a freight negotiation agent calling on behalf of {EMPRESA} to
negotiate the payment for a trip with a carrier. Your goal is to close a
deal that is fair to both parties, within financial limits that the system
controls — you do not know the absolute maximum limit; there is a function
that validates each offer before you say it.

TRIP DATA (injected by the system):
- Origin: {ORIGEN}
- Destination: {DESTINO}
- Type of freight / unit: {TIPO_CARGA}
- Internal target price (your reference, NEVER say it aloud): {PRECIO_OBJETIVO}
- Non-monetary compensations available to offer: {COMPENSACIONES_DISPONIBLES}
- Carrier's name (if known): {NOMBRE_TRANSPORTISTA}

ABSOLUTE RULE ON PRICES:
Never say an amount aloud without first having called the function
proponer_precio(monto) and having received "aceptado" (accepted). If the
function answers "rechazado" (rejected), do not propose that amount under any
circumstance — adjust downwards and continue. There is no scenario in which
you say a number that the function has not validated first.

Never reveal, hint or confirm that a maximum price or an upper limit exists.
If the carrier presses directly ("what is your ceiling?"), answer that the
amount offered reflects a fair market rate for this route — never say
phrases such as "I cannot pay more" or "that is my limit".

NEGOTIATION STRATEGY:
1. Open with one or two short sentences: greet, say that you are calling on
   behalf of {EMPRESA} and confirm that you are speaking with the right
   person. Do not make up a name for yourself or say a placeholder. Stop
   there and yield the turn — do not chain the route, the freight or the
   offer in that first message. Wait for the carrier to answer; the route
   and the type of freight come later, in a separate turn, and the first
   offer in another.
2. Make your first offer moderately below your internal target price —
   leave room to concede. (The exact anchoring percentage is a parameter to
   calibrate in Phase 5; in v0 use a prudent margin.)
3. If the carrier asks for more, do not concede on price at once: first
   offer a non-monetary compensation from the available list.
4. If the carrier insists on price after the compensation, raise your offer
   in decreasing increments — each increase must be smaller than the
   previous one. Validate each new offer with proponer_precio before saying it.
5. Never exceed your internal target price without first exhausting the
   available non-monetary compensations.
6. Before closing, repeat aloud the final amount and the agreed conditions
   in your own words, and ask for explicit confirmation (for example: "are
   we agreed, with that amount and those conditions?"). Only then call
   registrar_acuerdo(monto_final, condiciones).
7. If there is no agreement after three counter-offers, or the carrier
   refuses definitively, close politely and call
   finalizar_sin_acuerdo(motivo).

HANDLING OBJECTIONS:
- "The price is low for the distance/fuel": acknowledge the point, give
  market context, consider a non-monetary compensation before moving the price.
- "I need an advance payment": confirm the payment terms if you know them;
  if not, offer to follow up after the call.
- "I need paid waiting time at loading/unloading": it is a valid
  non-monetary compensation — offer it as part of the deal.
- Prolonged silence: wait a few seconds, ask naturally whether the carrier
  is still on the line; if it persists, close and call
  finalizar_sin_acuerdo("sin respuesta").
- Flat refusal or hang-up: do not insist or call again in the same session;
  close with finalizar_sin_acuerdo and the corresponding reason.

VOICE STYLE:
Neutral Spanish, short and natural sentences — this is synthesized as audio,
not read. Avoid lists and technical terms. Use natural acknowledgements
("entiendo", "tiene sentido", "claro") to sound human. Do not interrupt:
wait for the carrier to finish speaking. If the carrier answers in English,
switch to English naturally for the rest of the call.
\end{lstlisting}

\subsection{What each instruction is for}
\begin{itemize}[leftmargin=1.4cm]
  \item \textbf{Absolute rule on prices:} the maximum price is enforced in code,
  not in the text, so that the model can neither reveal it nor exceed it.
  \item \textbf{Not revealing the ceiling:} prevents the agent from giving away its
  negotiating position when asked directly, as in scenario~3.
  \item \textbf{Short opening that yields the turn:} lets the carrier come in;
  without the carrier's turns there is no latency to measure.
  \item \textbf{No made-up names or placeholders:} prevents the synthesizer from
  saying a filler text such as ``[name]''.
  \item \textbf{First offer below the target:} leaves room to concede during the
  negotiation.
  \item \textbf{Compensation before price, and decreasing concessions:} the
  negotiation rules of the business; scenarios~2 and~4 exercise them.
  \item \textbf{Confirmation before registering:} leaves a verifiable close.
  \item \textbf{Closing after three counter-offers:} sets a limit so that the call
  does not drag on without agreement; scenarios~2 and~3 can end this way.
  \item \textbf{Handling objections:} gives a prepared answer to the typical
  objections of the trade. The fuel objection appears in scenarios~2 and~3, and
  the waiting-time objection in scenario~4. Prolonged silence and hang-up define
  how to close if the carrier stops answering or cuts the call; no scenario of
  this study exercises them.
  \item \textbf{Voice style:} the text is turned into audio; lists and technical
  terms sound wrong when spoken.
  \item \textbf{Switch to English:} supports English with the same text. It was not
  evaluated: no script has lines in English (Section~4.1 of the article).
\end{itemize}

\subsection{Tools}

The four functions and their descriptions are the same in the three architectures.
Table~\ref{tab:herramientas} gives the exact text that the model receives, with its
translation.

\begin{table}[H]
\centering
\caption{The four tools: description and parameters as the model receives them
(in Spanish), and their translation.}
\label{tab:herramientas}
\scriptsize
\renewcommand{\arraystretch}{1.2}
\begin{tabularx}{\linewidth}{@{}P{2.7cm}YY@{}}
\toprule
\textbf{Function} & \textbf{Text received by the model} & \textbf{English translation}\\
\midrule
\texttt{proponer\_precio} & Valida un monto ANTES de pronunciarlo en voz alta. Responde "aceptado" o "rechazado". Si responde "rechazado", no digas ese monto: ajusta a la baja y vuelve a validar. Nunca menciones un monto que no haya sido aceptado por esta función. & Validates an amount BEFORE saying it aloud. Answers ``aceptado'' (accepted) or ``rechazado'' (rejected). If it answers ``rechazado'', do not say that amount: adjust downwards and validate again. Never mention an amount that this function has not accepted.\\
\quad\texttt{monto} (number) & Monto en pesos mexicanos que quieres ofrecer al transportista por el viaje completo. & Amount in Mexican pesos that you want to offer the carrier for the whole trip.\\
\addlinespace
\texttt{ofrecer\_\allowbreak compensacion\_\allowbreak no\_\allowbreak monetaria} & Registra que ofreces una compensación no monetaria del catálogo autorizado. Responde con la redacción exacta que debes usar al decirla en voz alta. Ofrece compensaciones antes de subir el precio. & Registers that you offer a non-monetary compensation from the authorized catalogue. Answers with the exact wording that you must use when saying it aloud. Offer compensations before raising the price.\\
\quad\texttt{tipo} (text) & Identificador de la compensación del catálogo. No inventes valores fuera de esta lista. Valores: \texttt{tiempo\_\allowbreak espera\_\allowbreak carga}, \texttt{prioridad\_\allowbreak proxima\_\allowbreak ruta}, \texttt{flexibilidad\_\allowbreak fecha\_\allowbreak entrega}. & Identifier of the compensation in the catalogue. Do not make up values outside this list. Values: waiting time at loading, priority for the next route, flexibility in the delivery date.\\
\addlinespace
\texttt{registrar\_\allowbreak acuerdo} & Cierra la negociación con acuerdo. Llámala solo después de haber repetido en voz alta el monto final y las condiciones, y de que el transportista lo haya confirmado explícitamente. & Closes the negotiation with an agreement. Call it only after having repeated aloud the final amount and the conditions, and after the carrier has confirmed them explicitly.\\
\quad\texttt{monto\_final} (number) & Monto acordado en pesos mexicanos. Debe ser uno que proponer\_precio ya haya aceptado. & Agreed amount in Mexican pesos. It must be one that proponer\_precio has already accepted.\\
\quad\texttt{condiciones} (list of texts) & Condiciones acordadas, en las palabras en que se pactaron: compensaciones incluidas, fechas, tiempos de espera. & Agreed conditions, in the words in which they were settled: compensations included, dates, waiting times.\\
\addlinespace
\texttt{finalizar\_\allowbreak sin\_\allowbreak acuerdo} & Cierra la negociación sin acuerdo, después de despedirte cordialmente. & Closes the negotiation without an agreement, after saying goodbye politely.\\
\quad\texttt{motivo} (text) & Por qué no hubo acuerdo: sin\_respuesta si dejó de contestar, rechazo\_definitivo si rechazó de forma tajante, sin\_acuerdo\_tras\_contraofertas si se agotaron las contraofertas, colgado si cortó la llamada, falla\_tecnica si hubo un problema de audio, otro en cualquier otro caso. & Why there was no agreement: no answer if the carrier stopped answering, definitive refusal if the carrier refused flatly, no agreement after counter-offers if the counter-offers ran out, hung up if the carrier cut the call, technical failure if there was an audio problem, other in any other case.\\
\bottomrule
\end{tabularx}
\end{table}

\paragraph{What each function answers.} \texttt{proponer\_precio} answers only
``aceptado'' or ``rechazado'', without saying why. It rejects the amount if it is not
a positive number, if it exceeds the maximum price (21,000 pesos) or if it breaks
the rule of decreasing concessions: when the amount is above the highest offer
already accepted, the increase over that offer must be strictly smaller than the
last accepted increase. Repeating or lowering an offer is always accepted. In any
other case it answers ``aceptado'' and records the amount in the history of offers.
\texttt{ofrecer\_compensacion\_no\_monetaria} answers with the exact wording of the
compensation (that of Table~\ref{tab:viaje}). The four functions run on the same
negotiation engine in the three architectures.


\section{Carrier lines and branching rules}
\label{sec:lineas}

The \RGuionLineas{} lines that the carrier can say are written in the Spanish of
Mexican carriers (``órale'', ``nomás'', ``parejo''), as a carrier would speak on the
phone. Each of the five scenarios has 8 turns and a different carrier profile, and
all five use the same trip data (Table~\ref{tab:viaje}). The negotiation is the
vehicle: the lines are written so that the agent has to use its tools in the middle
of the conversation, and to put into the audio what is hard for voice technology
(amounts, place names, trade terms).

Each scenario follows, with the reason for its lines and a table with all of them.
In the turns with branches, the column ``When it plays'' says what the agent must
have done for each line to play; ``---'' means that the turn has a single line.
The conditions of a turn are listed in the order in which the harness checks them.
Each line is given as the audio says it, which is what was played in the tests,
with an English translation; the asterisk marks the lines in which the audio does
not say exactly what the written script says (Section~\ref{sec:correcciones}).

\need{9cm}
\subsection{Scenario 1: cooperative}

This is the baseline conversation, against which the other four are read: a willing
carrier who asks about the freight, the date and the amount, hesitates once
(``Mmm'') and accepts. The lines are short and the carrier says no amounts, to
measure latency and recognition in the simplest case. The carrier's questions make
the agent give the trip data and its offer, and the last two lines close the deal,
so that the agent confirms it and registers it with \texttt{registrar\_acuerdo}.

\begin{scripttable}{Scenario 1, cooperative: the carrier's lines.}{tab:guion-e1}
T1 & --- & Sí, soy Carlos. Dígame. & Yes, this is Carlos. Go ahead.\\
T2 & --- & Va, suena interesante. ¿Me cuenta bien de la carga primero? & Okay, sounds interesting. Can you tell me about the load first?\\
T3 & --- & Ajá. ¿Y para cuándo la necesita en Laredo?\textsuperscript{*} & Uh-huh. And when do you need it in Laredo?\\
T4 & --- & Está bien. ¿Y en cuánto me dijo que quedaba el flete? & All right. And how much did you say the rate was?\\
T5 & --- & Mmm. Está un poquito bajo, pero puede funcionar. & Hmm. It's a little low, but it could work.\\
T6 & --- & ¿Sabe qué? Si me lo deja parejo, lo tomo. & You know what? If you make it a round number, I'll take it.\\
T7 & --- & Perfecto. Quedamos así entonces. & Perfect. We have a deal, then.\\
T8 & --- & Sí, así quedamos. Confirmado. & Yes, that's the deal. Confirmed.\\
\end{scripttable}

\need{15cm}
\subsection{Scenario 2: moderate haggler}

A carrier who asks for more again and again, with two objections that are common in
the trade: the price of diesel and how much the route pays. Each line pushes the
agent to go up, so the agent has to validate new amounts with
\texttt{proponer\_precio} several turns in a row. The amount requested appears in
two forms, in words (``dieciocho mil'') and in digits (``18,000''): they sound the
same in the audio, but the recognizer may write them differently. In turn~6 the
carrier self-corrects in mid-sentence (``Diecisiete\ldots{} no, dieciocho''), a hard
case for keeping the right amount. Three turns have branches, so that the
conversation keeps making sense whatever the agent answers: if the agent has
already reached 18,000, the carrier accepts; if it offered a compensation instead
of going up, the carrier insists on the number; otherwise the carrier keeps
pressing and, in the end, closes without an agreement.

\begin{scripttable}{Scenario 2, moderate haggler: the carrier's lines.}{tab:guion-e2}
T1 & --- & Sí, soy Carlos. A ver, dígame qué trae. & Yes, this is Carlos. Let's see, what have you got?\\
T2 & --- & No, está bajo para Monterrey–Laredo. El diesel está bien caro ahorita. & No, that's low for Monterrey–Laredo. Diesel is really expensive right now.\\
T3 & --- & ¿A cómo andan pagando esa ruta ustedes?\textsuperscript{*} & What are you paying for that route?\\
T4 & --- & Súbale tantito. Yo ando pensando en dieciocho mil. & Go up a little. I'm thinking eighteen thousand.\\
T5 & if the agent has already said 18{,}000 or more & Órale, así sí me late. & All right, now that works for me.\\
 & if the agent offered a compensation in its reply & Está bien lo de la compensación, pero el número tiene que ser dieciocho.\textsuperscript{*} & The compensation is fine, but the number has to be eighteen.\\
 & in any other case & Ya casi. Nada más déjelo en 18,000 parejo. & Almost there. Just make it an even 18,000.\\
T6 & --- & Ándele, es una buena ruta y se la hago rápida. Diecisiete… no, dieciocho.\textsuperscript{*} & Come on, it's a good route and I'll get it done fast. Seventeen… no, eighteen.\\
T7 & if the agent has already said 18{,}000 or more & Va. En dieciocho mil quedamos. & Okay. We're set at eighteen thousand.\\
 & in any other case & ¿Seguro que no le puede mover más? & Are you sure you can't move it any higher?\\
T8 & if the agent has already said 18{,}000 or more & Sale, quedamos en dieciocho mil. Confirmado. & Deal, we're set at eighteen thousand. Confirmed.\\
 & in any other case & No, en menos no me sale. Ahí déjelo, contácteme si se puede. & No, for less it doesn't work for me. Leave it there; contact me if it becomes possible.\\
\end{scripttable}

\need{10.5cm}
\subsection{Scenario 3: aggressive, with a high anchor}

A curt carrier who opens far above the maximum price: the carrier asks for 28,000,
then 25,000, and does not give way; the maximum is 21,000 (Table~\ref{tab:viaje}).
It exercises the limit enforced in code: \texttt{proponer\_precio} would reject any
attempt by the agent to approach what the carrier asks. In turn~5 the carrier asks
directly for the ceiling (``¿Cuánto es lo máximo que puede pagar?''), the situation
in which the instructions forbid the agent to reveal it. The high figures are
repeated in words and in digits, and the call ends without an agreement (``Entonces
no hay trato''), so that the agent closes with \texttt{finalizar\_sin\_acuerdo}.

\begin{scripttable}{Scenario 3, aggressive with a high anchor: the carrier's lines.}{tab:guion-e3}
T1 & --- & Bueno, soy Carlos. ¿De qué se trata? & Hello, this is Carlos. What's this about?\\
T2 & --- & Uy no, para nada. Ese viaje me lo pagan a veintiocho mil. & Oh no, no way. I get paid twenty-eight thousand for that trip.\\
T3 & --- & Es que ustedes no saben cómo está el diesel y las casetas. 28,000, mínimo. & You people don't know what diesel and tolls cost. 28,000, minimum.\\
T4 & --- & No le hace que le parezca caro. 28,000 o no hay viaje. & I don't care if it seems expensive to you. 28,000 or there's no trip.\\
T5 & --- & A ver, dígame cuál es su tope. ¿Cuánto es lo máximo que puede pagar? & Okay, tell me what your cap is. What's the most you can pay?\\
T6 & --- & Bueno… le bajo a veinticinco mil, pero de ahí…\textsuperscript{*} & Fine… I'll come down to twenty-five thousand, but from there…\\
T7 & --- & No, veinticinco es mi último precio. Ni un peso menos. & No, twenty-five is my final price. Not one peso less.\\
T8 & --- & Entonces no hay trato. Ahí muere. & Then there's no deal. That's the end of it.\\
\end{scripttable}

\need{13.5cm}
\subsection{Scenario 4: non-monetary objection}

A carrier who finds the price acceptable but sets conditions: the time spent
waiting at loading and the flexibility in the delivery date. The carrier says no
amounts, so that the agent has to understand the objection and answer with a
compensation (\texttt{ofrecer\_compensacion\_no\_monetaria}) instead of raising the
price. For this reason three turns have branches: if the agent offers a
compensation, the carrier moves forward (``Ah, eso sí me sirve''); if the agent
answers with money, the carrier objects (``No le pedí más dinero''). The scenario
has long questions and the vocabulary of conditions (``tiempo de carga'', ``fecha
de entrega'', ``por escrito'').

\begin{scripttable}{Scenario 4, non-monetary objection: the carrier's lines.}{tab:guion-e4}
T1 & --- & Sí, soy Carlos. Dígame. & Yes, this is Carlos. Go ahead.\\
T2 & --- & Mire, el precio ahí más o menos. Pero tengo otras broncas con esa ruta. & Look, the price is more or less okay. But I have other problems with that route.\\
T3 & --- & El problema es el tiempo de carga. Siempre me tienen horas esperando. & The problem is the loading time. They always keep me waiting for hours.\\
T4 & if the agent offered a compensation in its reply & Ah, eso sí me sirve. ¿Y me lo garantizan por escrito? & Ah, that does help me. And will you guarantee it in writing?\\
 & in any other case & No le pedí más dinero. Le pregunté por el tiempo de espera. & I didn't ask you for more money. I asked you about the waiting time.\\
T5 & --- & ¿Y qué tan flexibles son con la fecha de entrega? A veces se me empalma. & And how flexible are you with the delivery date? Sometimes my jobs overlap.\\
T6 & if the agent offered a compensation in its reply & Órale, con eso ya cambia la cosa. & All right, that changes things.\\
 & in any other case & Otra vez el precio. Le estoy hablando de condiciones, no de dinero. & The price again. I'm talking to you about conditions, not money.\\
T7 & --- & Bueno. Con las condiciones que me da, ¿en cuánto queda el flete al final? & Okay. With the conditions you're giving me, what does the rate come to in the end?\\
T8 & if the agent offered a compensation at any point in the call & Va, quedamos con eso. Confirmado. & Okay, we're set with that. Confirmed.\\
 & in any other case & Pues sin arreglar lo del tiempo, no me late. Déjelo pendiente. & Well, without sorting out the waiting time, I don't like it. Let's leave it pending.\\
\end{scripttable}

\need{11.5cm}
\subsection{Scenario 5: recognition stress test}

A cooperative carrier whose challenge lies not in the negotiation but in what is
said: each line packs what is hardest for a recognizer on telephone audio. There
are numbers that are not money (``caja seca de 53 pies'', ``tara de catorce
toneladas'', ``veintidós de carga útil''), a long amount in words (``mil ciento
veinte pesos''), a road (``la 85''), two place names (Nuevo Laredo and Sabinas
Hidalgo), trade terms (\emph{tara}, \emph{carta porte}, \emph{casetas}) and a
self-correction of an amount (``dieciocho mil\ldots{} no, dieciocho mil
quinientos'') which is later repeated in digits (``18,500''). The call is normal
and ends in an agreement, like the others: the difficulty is in the words, not in
the carrier's behaviour.

\begin{scripttable}{Scenario 5, recognition stress test: the carrier's lines.}{tab:guion-e5}
T1 & --- & Sí, soy Carlos, mande. & Yes, this is Carlos, go ahead.\\
T2 & --- & Déjeme apuntar. Traigo caja seca de 53 pies, con tara de catorce toneladas y me da veintidós de carga útil. & Let me write this down. I run a 53-foot dry van, with a tare weight of fourteen tons, and it gives me twenty-two of payload.\\
T3 & --- & Es que de aquí a Nuevo Laredo me voy por la 85, y nomás en casetas se me van mil ciento veinte pesos. & From here to Nuevo Laredo I take Highway 85, and the tolls alone cost me one thousand one hundred twenty pesos.\\
T4 & --- & Mire, no me siento a gusto con ese flete. Yo andaba pensando en dieciocho mil… no, dieciocho mil quinientos. & Look, I'm not comfortable with that rate. I was thinking eighteen thousand… no, eighteen thousand five hundred.\\
T5 & --- & Otra cosa, la carta porte la necesito bien hecha. El mes pasado un flete tuvo que esperar en Sabinas Hidalgo por un dato mal puesto. & Another thing, I need the carta porte (the bill of lading) done right. Last month a load had to wait in Sabinas Hidalgo because of one wrong entry.\\
T6 & --- & Ándele, con eso ya me late. Ya con la tara del remolque, menos no me sale. & Come on, now that works for me. With the trailer's tare weight, I can't do it for less.\\
T7 & --- & Órale, pues vaya que sí. Quedamos en 18,500.\textsuperscript{*} & Right on, sure thing. We're set at 18,500.\\
T8 & --- & Sí, así quedamos, confirmado. Mándeme la carta porte al correo. & Yes, that's the deal, confirmed. Send me the carta porte by email.\\
\end{scripttable}

\need{12.5cm}
\subsection{Interruption scenario}

These are twelve short lines that the harness plays over the agent while it is
speaking (Section~2.7 of the article). They are generic on purpose: they fall on
whatever response is in progress, so they have to make sense at any point of the
negotiation. There are two kinds, each with what is expected of the agent. In the
six interruptions the carrier takes the floor with a question or an objection
(``Oiga, ¿eso incluye las casetas?''): the agent should stop and answer. In the six
backchannels (``ajá'', ``mjm'', ``sí'') the carrier only signals that it is still
listening: the agent should keep talking.

\begin{table}[H]
\centering
\caption{Interruption scenario: the lines that the harness plays over the agent.}
\label{tab:guion-int}
\footnotesize
\renewcommand{\arraystretch}{1.15}
\setlength{\tabcolsep}{4pt}
\begin{tabular}{@{}P{2.1cm}P{6.4cm}P{\dimexpr\linewidth-8.5cm-4\tabcolsep\relax}@{}}
\toprule
\textbf{Kind} & \textbf{Line (what the audio says)} & \textbf{English translation}\\
\midrule
Interruption & Espéreme, espéreme… ¿cuánto dijo? & Hold on, hold on… how much did you say?\\
 & No, no, perdón, déjeme decirle algo. & No, no, sorry, let me tell you something.\\
 & Oiga, ¿eso incluye las casetas? & Hey, does that include the tolls?\\
 & Ya, ya le entendí. ¿Y el anticipo? & Okay, okay, I got it. And the advance payment?\\
 & Momento, eso no fue lo que quedamos. & Wait a moment, that's not what we agreed.\\
 & Perdón que lo interrumpa, pero así no me sale. & Sorry to interrupt, but that doesn't work for me.\\
Backchannel & Ajá. & Uh-huh.\\
 & Mjm. & Mm-hm.\\
 & Sí. & Yes.\\
 & Ok. & OK.\\
 & Ándale. & Right.\\
 & Sí, sí. & Yeah, yeah.\\
\bottomrule
\end{tabular}
\end{table}

\subsection{How the harness chooses the line of each branch}
\label{sec:ramas}

In \RGuionTurnosConRamas{} of the \RGuionBase{} turns (\RGuionTurnosConRamasLista,
where E is the scenario and T the turn) the line that plays depends on what the
agent said in its previous response: these are the branches, which add
\RGuionAlt{} alternative lines to the \RGuionBase{} base lines. The conditions are
in the column ``When it plays'' of the tables above. The harness applies them as
follows:
\begin{itemize}[leftmargin=1.4cm]
  \item The branches of a turn are checked in order, and the first whose condition
  holds wins. The last one has no condition: it plays if none of the previous ones
  held.
  \item The decision rests on what the agent \emph{said} (the text of its
  response), not on the functions it called, as a carrier would decide.
  \item \textbf{Amounts:} they are read with the same extractor as the error on
  amounts (Section~2.6 of the article), so ``dieciocho mil'' and ``18,000'' are the
  same number; only quantities of 1,000 or more count. The condition ``has already
  said 18,000 or more'' compares against the highest amount that the agent said in
  the whole call.
  \item \textbf{Compensations:} they are recognized by fixed phrases. Waiting time:
  ``tiempo de espera'', ``espera pagado'', ``espera pagada'', ``horas de espera''.
  Priority: ``prioridad'', ``siguiente ruta'', ``próxima ruta''. Date flexibility:
  ``día adicional'', ``día más de flexibilidad'', ``flexibilidad en la fecha'',
  ``flexibilidad en la entrega''. A phrase does not count if one of the three
  preceding words is a negation (\emph{no}, \emph{sin}, \emph{tampoco},
  \emph{ninguna}, \emph{ningún}, \emph{ninguno}, \emph{nunca}, \emph{imposible}),
  so that ``no tenemos flexibilidad en la fecha'' (``we have no flexibility in the
  date'') is not taken as an offer.
\end{itemize}

\subsection{Lines in which the audio differs from the written script}
\label{sec:correcciones}

In \RGuionCorregidas{} of the \RGuionLineas{} lines the audio does not say exactly
what the written script says: \RGuionCorregidasRecorte{} because trimming the
silence removed the end of the line, and \RGuionCorregidasDesvio{} because the
recording departed from the script. The tables above give what the audio says and
mark these lines with an asterisk; the reference for the word error rate also uses
the audio (Section~2.6 of the article). For each of these lines, identified by the
name of its audio file (scenario, turn and, where applicable, branch), the written
script said:
{\footnotesize
\begin{itemize}[leftmargin=1.2cm,itemsep=0pt]
\item \texttt{e5\_t7}: the written script said «Órale, pues vaya que sí. Quedamos en 18,500 con todo incluido.»
\item \texttt{e3\_t6}: the written script said «Bueno… le bajo a veinticinco mil, pero de ahí no me muevo.»
\item \texttt{e2\_t6}: the written script said «Ándele, es una buena ruta y se la hago rápida. Diecisiete… no, dieciocho mil.»
\item \texttt{e2\_t5\_alt1}: the written script said «Está bien lo de la compensación, pero el número tiene que ser dieciocho mil.»
\item \texttt{e1\_t3}: the written script said «Ajá. ¿Y para cuándo la necesitan en Laredo?»
\item \texttt{e2\_t3}: the written script said «¿A cómo andan pagando esa ruta ustedes normalmente?»
\end{itemize}
}


\section{Hypotheses tested in the interruption scenario}
\label{sec:hipotesis}

Eight hypotheses were set before the interruption campaign (Sections~2.7 and~2.9 of
the article). Each one compares one value per call between two conditions, paired
by scenario and repetition, with the two-sided Wilcoxon signed-rank test. With
Bonferroni correction over the \RIntPruebaFamilia{} tests, a hypothesis passes if
$p<0.05/\RIntPruebaFamilia=\RIntPruebaAlfa$. Table~\ref{tab:hipotesis} lists all of
them. The number of pairs is lower than the \RIntLlamadasA{} calls per architecture
when a measure is not defined in every call: the pause exists only in calls in
which the agent stopped at a backchannel and spoke again within 8~s, and the
fraction of resumptions that continued the sentence only in calls with at least one
classified resumption.

\begin{table}[H]
\centering
\caption{The eight hypotheses of the interruption scenario. One value per call;
pairs by scenario and repetition; two-sided Wilcoxon signed-rank test. A hypothesis
passes if $p<\RIntPruebaAlfa$ (Bonferroni over \RIntPruebaFamilia{} tests).}
\label{tab:hipotesis}
\footnotesize
\renewcommand{\arraystretch}{1.2}
\begin{tabularx}{\linewidth}{@{}rYP{3.5cm}ccc@{}}
\toprule
\textbf{\#} & \textbf{Measure (one value per call)} & \textbf{Comparison} & \textbf{Pairs} & \textbf{$p$} & \textbf{Passes}\\
\midrule
1 & Fraction of backchannels at which the agent stopped & GPT-Live vs Pipecat & \RIntPruebaCalloTresUnoN & \RIntPruebaCalloTresUnoP & \RIntPruebaCalloTresUnoPasa\\
2 & Fraction of backchannels at which the agent stopped & GPT-Live vs LiveKit & \RIntPruebaCalloTresDosN & \RIntPruebaCalloTresDosP & \RIntPruebaCalloTresDosPasa\\
\addlinespace
3 & Median duration of the pause after stopping at a backchannel & GPT-Live vs Pipecat & \RIntPruebaPausaTresUnoN & \RIntPruebaPausaTresUnoP & \RIntPruebaPausaTresUnoPasa\\
4 & Median duration of the pause after stopping at a backchannel & GPT-Live vs LiveKit & \RIntPruebaPausaTresDosN & \RIntPruebaPausaTresDosP & \RIntPruebaPausaTresDosPasa\\
5 & Median duration of the pause after stopping at a backchannel & Pipecat vs LiveKit & \RIntPruebaPausaUnoDosN & \RIntPruebaPausaUnoDosP & \RIntPruebaPausaUnoDosPasa\\
\addlinespace
6 & Fraction of events at which the agent stopped & GPT-Live: backchannels vs controls & \RIntPruebaApoyoControlTresN & \RIntPruebaApoyoControlTresP & \RIntPruebaApoyoControlTresPasa\\
\addlinespace
7 & Fraction of resumptions that continued the sentence where it had stopped & GPT-Live vs Pipecat & \RIntPruebaSigueTresUnoN & \RIntPruebaSigueTresUnoP & \RIntPruebaSigueTresUnoPasa\\
8 & Fraction of resumptions that continued the sentence where it had stopped & GPT-Live vs LiveKit & \RIntPruebaSigueTresDosN & \RIntPruebaSigueTresDosP & \RIntPruebaSigueTresDosPasa\\
\bottomrule
\end{tabularx}

\raggedright\footnotesize The agent stopped if its voice was cut by a silence of at
least 0.7~s starting within 3~s of the event. Controls are the instants at which
the harness plays nothing. Resumptions classified as ``cannot tell'' do not count
in hypotheses 7 and~8.
\end{table}

\end{document}
